\section{Revisão sistemática da literatura (PRISMA 2020)}
\label{sec:revisao}

A atividade~1 do cronograma (revisão bibliográfica sobre MCS, DCS/DES,
diversidade e Oracle) foi conduzida como revisão sistemática segundo o PRISMA
2020 (Page et al., 2021). Protocolo, strings, exportações brutas, decisões por
registro e scripts estão em \texttt{docs/revisao-sistematica/}; as contagens do
fluxograma são recalculadas por script a partir dos CSVs.

\paragraph{Perguntas.} RQ1 --- como o Oracle é definido e usado como limite
superior em MCS e DCS/DES? RQ2 --- que trabalhos propõem limites mais realistas
ou generalizações do Oracle? RQ3 --- como diversidade e redundância de acertos se
relacionam com a folga entre o Oracle e os combinadores reais? RQ4 --- que
protocolos e bases são usados?

\paragraph{Elegibilidade.} Publicado entre $2021$ e $2026$; em inglês, português
ou chinês; trata do Oracle ou de limite superior em MCS/DCS/DES, \emph{ou} mede
diversidade $\times$ desempenho do ensemble empiricamente. Exclusão: preprints,
``oracle'' em outro sentido (blockchain, oráculo de teste, banco de dados),
texto inacessível, duplicatas. Os trabalhos fundadores anteriores a $2021$
(Kuncheva e Whitaker, 2003; Souza e Cavalcanti, 2017; Cruz et al., 2017, 2018)
ficam fora do corpus por construção e aparecem nas demais seções como
referencial teórico.

\paragraph{Fontes e buscas.} Scopus (pelo Portal de Periódicos CAPES, acesso
CAFe), IEEE Xplore e OpenAlex, com quatro strings --- S1 Oracle e limites, S2
diversidade $\times$ acurácia, S3 seleção dinâmica, S4 geração de pool --- e
snowballing para frente pelo campo \texttt{REF} do Scopus (quem cita Monteiro et
al., 2023, e Souza e Cavalcanti, 2017). O Web of Science estava planejado e não
foi consultado: toda busca devolve erro pelo proxy da CAPES. ScienceDirect e
SpringerLink entram pela indexação do Scopus e serviram ao acesso ao texto
completo. Buscas executadas em $23$ de outubro de $2026$.

\begin{figure}[htbp]
\centering
\begin{tikzpicture}[
  node distance=13mm and 8mm,
  box/.style={draw, rounded corners=2pt, align=left, font=\footnotesize,
              text width=6.2cm, inner sep=4pt},
  side/.style={draw, rounded corners=2pt, align=left, font=\footnotesize,
               text width=5.6cm, inner sep=4pt, fill=black!4},
  arr/.style={-{Stealth[length=2mm]}}]
\node[box] (id) {\textbf{Identificação:} $677$ registros\\
  Scopus $319$ $\cdot$ IEEE Xplore $108$ $\cdot$ OpenAlex $250$};
\node[side, right=of id] (dup) {Duplicatas removidas: $329$};
\node[box, below=of id] (scr) {\textbf{Triagem} (título e resumo): $348$};
\node[side, right=of scr] (exs) {Excluídos: $236$\\
  aplicação sem análise $148$; fora do tema ou ``oracle'' em outro sentido
  $42$; não-classificação $31$; não revisado por pares $10$; sem informação $5$};
\node[box, below=of scr] (sou) {Buscados em texto completo: $112$};
\node[side, right=of sou] (nr) {Não recuperados: $10$};
\node[box, below=of sou] (ass) {\textbf{Elegibilidade} (texto completo): $102$};
\node[side, right=of ass] (exf) {Excluídos: $78$\\
  não usa o Oracle nem relaciona diversidade e acurácia $74$; ``oracle'' em
  outro sentido $2$; retratado ou outro relato $2$};
\node[box, below=of ass] (inc) {\textbf{Incluídos:} $24$ relatos, $22$ estudos\\
  generaliza o Oracle $1$ $\cdot$ Oracle como referência $10$ $\cdot$
  diversidade $\times$ acurácia $13$};
\draw[arr] (id) -- (scr); \draw[arr] (scr) -- (sou);
\draw[arr] (sou) -- (ass); \draw[arr] (ass) -- (inc);
\draw[arr] (id) -- (dup); \draw[arr] (scr) -- (exs);
\draw[arr] (sou) -- (nr); \draw[arr] (ass) -- (exf);
\end{tikzpicture}
\caption{Fluxograma PRISMA 2020 da revisão. As contagens saem de
\texttt{fluxo\_prisma.json} e fecham por construção: $677-329=348$,
$348-236=112$, $102-78=24$.}
\label{fig:prisma}
\end{figure}

\subsection{Síntese por pergunta}

\paragraph{RQ1 --- o Oracle é reportado, não analisado.} Dez estudos usam o
Oracle como referência experimental de teto, quase sempre como uma coluna a mais
na tabela de resultados (FH-DES, MLS, GNN-DES, Monteiro et al., entre outros). Em
parte deles o Oracle é retirado dos testes estatísticos por ser ``uma
abstração''; só um deixa a folga visível por base (MLS: Balance-Scale $0{,}910$
contra $0{,}987$ do Oracle). Duas linhas o usam como rótulo de meta-aprendizado
(o seletor é treinado para imitar o Oracle no treino), e um estudo o apresenta
como ``meta-classificador'' implantável --- o que é vazamento de rótulo, já que o
Oracle consulta a classe verdadeira. Com pools de $100$ classificadores o Oracle
fica perto de $1$ (Monteiro et al.: $\approx99\%$ para quase todo gerador), o que
confirma, fora deste projeto, que \On{1} pouco discrimina entre pools
(Seção~\ref{sec:obj1}).

\paragraph{RQ2 --- uma única generalização, com outro eixo.} O
\emph{Hardness-aware Oracle} (Cordeiro, Cavalcanti e Cruz, 2023) pondera o acerto pela
dificuldade da instância para escolher, a posteriori, entre métodos de DES.
Nenhum trabalho do período propõe níveis ``pelo menos $N$ de $M$ corretos'', uma
curva de Oracles ou a ligação $\On{M/2+1}\le\mvr\le\On{M/2}$. O \On{N} preenche
uma lacuna documentada, e é complementar ao HaO: um conta a redundância de
acertos, o outro pesa a dificuldade da instância.

\paragraph{RQ3 --- diversidade medida contra a acurácia, nunca contra a folga.}
Doze estudos relacionam diversidade e acurácia, e a maioria propõe uma medida
nova (\emph{double-fault disagreement}, diversidade de interpretabilidade,
distância de ranking, volume em embeddings aprendidos, ambiguidade). A relação
encontrada é fraca e depende do contexto: correlação positiva fraca em
sensoriamento remoto, sinal que muda com a taxa de amostragem, efeito
condicionado ao combinador em ensembles heterogêneos. Nenhum estudo relaciona
diversidade ou redundância com a fração do teto que um combinador real recupera
--- que é onde este projeto encontra sinal forte ($\dfe$, Spearman $-0{,}89$ com
a folga, Seção~\ref{sec:obj8}).

\paragraph{RQ4 --- protocolos.} O padrão é $20$--$35$ bases UCI/KEEL, pools de
$100$ Perceptrons (o mesmo do PGDCS), holdout replicado ou $k$-fold, e Wilcoxon,
Friedman ou Quade com o Oracle fora do teste. Imagens aparecem só via atributos
profundos (CNNs em CIFAR, ViT/CNN em texturas), o que respalda a bateria do
ADR~0022 (Seção~\ref{sec:obj5}). Nenhum estudo discute os dois vieses que este
projeto mediu: escolher o melhor método no teste e medir no mesmo conjunto que
seleciona.

\paragraph{Limitações da revisão.} A triagem foi feita por um revisor assistido
por IA, com critérios idênticos por lote e sem segundo revisor independente (sem
kappa). A elegibilidade em texto completo usou os trechos ao redor de cada
ocorrência de ``oracle'' e contagens de termos, com recheck integral dos
artigos excluídos que tinham muitas ocorrências. O Web of Science, as bases
chinesas (CNKI, Wanfang) e o snowballing para trás ficaram de fora. Recomenda-se
que o orientador revise uma amostra das exclusões. Os $24$ relatos incluídos
estão listados ao final, em ``Estudos incluídos na revisão sistemática''.
